\documentclass[10pt,a4paper]{amsart}
\usepackage{amsmath,amssymb,amsthm,graphicx}
% Retained minimal preamble for the Stacks Project chapter flat.tex (More on Flatness).
% Verbatim excerpt of the chapter preamble: Xy-pic setup (native preamble lines 14-18) and the native
% theorem declarations with the mathematical macros (native preamble lines 154-217). Copied mechanically;
% nothing added, removed or redefined. The wrapper class and the amsmath/amssymb/amsthm/graphicx base
% packages are supplied by the assembler, so the class/externaldocument part of the native preamble is omitted.
% For commutative diagrams we use Xy-pic
\usepackage[all]{xy}

% We use 2cell for 2-commutative diagrams.
\xyoption{2cell}
% Theorem environments.
%
\theoremstyle{plain}
\newtheorem{theorem}[subsection]{Theorem}
\newtheorem{proposition}[subsection]{Proposition}
\newtheorem{lemma}[subsection]{Lemma}

\theoremstyle{definition}
\newtheorem{definition}[subsection]{Definition}
\newtheorem{example}[subsection]{Example}
\newtheorem{exercise}[subsection]{Exercise}
\newtheorem{situation}[subsection]{Situation}

\theoremstyle{remark}
\newtheorem{remark}[subsection]{Remark}
\newtheorem{remarks}[subsection]{Remarks}

\numberwithin{equation}{subsection}

% Macros
%
\def\lim{\mathop{\mathrm{lim}}\nolimits}
\def\colim{\mathop{\mathrm{colim}}\nolimits}
\def\Spec{\mathop{\mathrm{Spec}}}
\def\Hom{\mathop{\mathrm{Hom}}\nolimits}
\def\Ext{\mathop{\mathrm{Ext}}\nolimits}
\def\SheafHom{\mathop{\mathcal{H}\!\mathit{om}}\nolimits}
\def\SheafExt{\mathop{\mathcal{E}\!\mathit{xt}}\nolimits}
\def\Sch{\mathit{Sch}}
\def\Mor{\mathop{\mathrm{Mor}}\nolimits}
\def\Ob{\mathop{\mathrm{Ob}}\nolimits}
\def\Sh{\mathop{\mathit{Sh}}\nolimits}
\def\NL{\mathop{N\!L}\nolimits}
\def\CH{\mathop{\mathrm{CH}}\nolimits}
\def\proetale{{pro\text{-}\acute{e}tale}}
\def\etale{{\acute{e}tale}}
\def\QCoh{\mathit{QCoh}}
\def\Ker{\mathop{\mathrm{Ker}}}
\def\Im{\mathop{\mathrm{Im}}}
\def\Coker{\mathop{\mathrm{Coker}}}
\def\Coim{\mathop{\mathrm{Coim}}}

% Boxtimes
%
\DeclareMathSymbol{\boxtimes}{\mathbin}{AMSa}{"02}

%
% Macros for moduli stacks/spaces
%
\def\QCohstack{\mathcal{QC}\!\mathit{oh}}
\def\Cohstack{\mathcal{C}\!\mathit{oh}}
\def\Spacesstack{\mathcal{S}\!\mathit{paces}}
\def\Quotfunctor{\mathrm{Quot}}
\def\Hilbfunctor{\mathrm{Hilb}}
\def\Curvesstack{\mathcal{C}\!\mathit{urves}}
\def\Polarizedstack{\mathcal{P}\!\mathit{olarized}}
\def\Complexesstack{\mathcal{C}\!\mathit{omplexes}}
% \Pic is the operator that assigns to X its picard group, usage \Pic(X)
% \Picardstack_{X/B} denotes the Picard stack of X over B
% \Picardfunctor_{X/B} denotes the Picard functor of X over B
\def\Pic{\mathop{\mathrm{Pic}}\nolimits}
\def\Picardstack{\mathcal{P}\!\mathit{ic}}
\def\Picardfunctor{\mathrm{Pic}}
\def\Deformationcategory{\mathcal{D}\!\mathit{ef}}

% --- appended 2026-10-12 (Chief): the project own definitions of the reference/slogan/history
% display environments, copied mechanically from the shared Stacks preamble lines 8-12
% (\usepackage{verbatim} plus the three \newenvironment lines, which the chapter preamble keeps
% verbatim from the project preamble). Nothing added, removed or redefined: these environments
% carry no mathematical content and the project itself typesets them as comments. Without these
% three lines a counted statement whose own text uses them fails with Environment slogan undefined.
\usepackage{verbatim}
\newenvironment{reference}{\comment}{\endcomment}
\newenvironment{slogan}{\comment}{\endcomment}
\newenvironment{history}{\comment}{\endcomment}
\usepackage{cleveref}
\crefname{originalsection}{original section}{original sections}
\Crefname{originalsection}{Original section}{Original sections}
\makeatletter
\newlabel{algebra-section-going-up}{{00HU}{source}{}{sourcepointer-1}{}}
\newlabel{algebra-section-going-up@cref}{{[originalsection][0][]00HU}{[1][0][]source}{}{}{}}
\makeatother
\crefname{originallemma}{original lemma}{original lemmas}
\Crefname{originallemma}{Original lemma}{Original lemmas}
\makeatletter
\newlabel{divisors-lemma-base-change-relative-assassin}{{05DC}{source}{}{sourcepointer-2}{}}
\newlabel{divisors-lemma-base-change-relative-assassin@cref}{{[originallemma][0][]05DC}{[1][0][]source}{}{}{}}
\makeatother
\crefname{originallemma}{original lemma}{original lemmas}
\Crefname{originallemma}{Original lemma}{Original lemmas}
\makeatletter
\newlabel{morphisms-lemma-base-change-flat}{{01U9}{source}{}{sourcepointer-3}{}}
\newlabel{morphisms-lemma-base-change-flat@cref}{{[originallemma][0][]01U9}{[1][0][]source}{}{}{}}
\makeatother
\crefname{originallemma}{original lemma}{original lemmas}
\Crefname{originallemma}{Original lemma}{Original lemmas}
\makeatletter
\newlabel{schemes-lemma-points-fibre-product}{{01JT}{source}{}{sourcepointer-4}{}}
\newlabel{schemes-lemma-points-fibre-product@cref}{{[originallemma][0][]01JT}{[1][0][]source}{}{}{}}
\makeatother
\title{Purity descends along a morphism flat at the point}
\author{The Stacks Project Authors}
\date{stacks-project:flat.tex:local-snapshot-20261011; GNU FDL 1.2 or later}
\begin{document}
\maketitle
\noindent\textit{Editorial note.} All delivered mathematical text is the author's own, copied from the pinned Stacks Project chapter file flat.tex; no statement or proof step was rewritten, reordered or removed.\par\medskip
\noindent\textit{Editorial note.} Preamble provenance: the chapter's own retained minimal preamble (sources/stacks-flat/retained-preamble.tex, lines 1-85) is a verbatim excerpt of the native Xy-pic setup and of the native theorem declarations and mathematical macros, followed by the project's own reference/slogan/history display-environment definitions copied mechanically from the shared preamble; the wrapper class and the amsmath/amssymb/amsthm/graphicx base packages are supplied by the assembler, and extra\_packages is left empty because everything needed is already in the retained preamble. No macro was invented or redefined.\par\medskip
\noindent\textit{Editorial note.} Same-chapter pointers are located at their \textbackslash{}label lines inside the pinned chapter file; cross-chapter pointers are located in the project's own tag list staged in the snapshot (sources/stacks-chow/tags.txt), at the verbatim lines carrying their canonical tags.\par\medskip
\noindent\textit{Allowed standard prerequisites (editorial).}\begin{itemize}
\item Flatness of the base change of a morphism flat at a point, and lifting of specialisations along a flat morphism. Morphisms, label morphisms-lemma-base-change-flat (proof line 4967); Algebra, Section algebra-section-going-up (proof line 4970). The first makes the projection from the base change flat at the lifted point, and the second produces a specialisation upstairs lying over the given one.
\item Points of a fibre product and base change of the relative assassin. Schemes, label schemes-lemma-points-fibre-product (proof line 4973); Divisors, label divisors-lemma-base-change-relative-assassin (proof line 4975). The first produces the required point of the fibre over the lifted specialisation and the second shows it is a relative associated point of the pulled back module.
\end{itemize}\medskip
\ifdefined\hypertarget\hypertarget{sourcepointer-1}{}\fi
\noindent\textit{Source reference.} Original section 00HU: sources/stacks-chow/tags.txt, source lines 656--656. Navigation pointer for \textbackslash{}ref\{algebra-section-going-up\} used inside the counted span (line 4970); the argument that consumes it is retained in the span.\par\medskip
\ifdefined\hypertarget\hypertarget{sourcepointer-2}{}\fi
\noindent\textit{Source reference.} Original lemma 05DC: sources/stacks-chow/tags.txt, source lines 6950--6950. Navigation pointer for \textbackslash{}ref\{divisors-lemma-base-change-relative-assassin\} used inside the counted span (line 4975); the argument that consumes it is retained in the span.\par\medskip
\ifdefined\hypertarget\hypertarget{sourcepointer-3}{}\fi
\noindent\textit{Source reference.} Original lemma 01U9: sources/stacks-chow/tags.txt, source lines 2402--2402. Navigation pointer for \textbackslash{}ref\{morphisms-lemma-base-change-flat\} used inside the counted span (line 4967); the argument that consumes it is retained in the span.\par\medskip
\ifdefined\hypertarget\hypertarget{sourcepointer-4}{}\fi
\noindent\textit{Source reference.} Original lemma 01JT: sources/stacks-chow/tags.txt, source lines 2026--2026. Navigation pointer for \textbackslash{}ref\{schemes-lemma-points-fibre-product\} used inside the counted span (line 4973); the argument that consumes it is retained in the span.\par\medskip
\begin{lemma}
\label{lemma-flat-descend-pure}
Let $f : X \to S$ be a morphism of schemes which is of finite type.
Let $\mathcal{F}$ be a finite type quasi-coherent $\mathcal{O}_X$-module.
Let $s \in S$. Let $(S', s') \to (S, s)$ be a morphism of pointed schemes.
Assume $S' \to S$ is flat at $s'$.
\begin{enumerate}
\item If $\mathcal{F}_{S'}$ is pure along $X_{s'}$,
then $\mathcal{F}$ is pure along $X_s$.
\item If $\mathcal{F}_{S'}$ is universally pure along $X_{s'}$,
then $\mathcal{F}$ is universally pure along $X_s$.
\end{enumerate}
\end{lemma}

\begin{proof}
Let $(T \to S, t' \leadsto t, \xi)$ be an impurity of
$\mathcal{F}$ above $s$. Set $T_1 = T \times_S S'$, and let $t_1$
be the unique point of $T_1$ mapping to $t$ and $s'$. Since
$T_1 \to T$ is flat at $t_1$, see
Morphisms, Lemma \ref{morphisms-lemma-base-change-flat},
there exists a specialization $t'_1 \leadsto t_1$ lying over
$t' \leadsto t$, see
Algebra, Section \ref{algebra-section-going-up}.
Choose a point $\xi_1 \in X_{t'_1}$ which corresponds to a generic
point of $\Spec(\kappa(t'_1) \otimes_{\kappa(t')} \kappa(\xi))$, see
Schemes, Lemma \ref{schemes-lemma-points-fibre-product}.
By
Divisors, Lemma \ref{divisors-lemma-base-change-relative-assassin}
we see that $\xi_1 \in \text{Ass}_{X_{T_1}/T_1}(\mathcal{F}_{T_1})$.
As the Zariski closure of $\{\xi_1\}$ in $X_{T_1}$ maps into the
Zariski closure of $\{\xi\}$ in $X_T$ we conclude that
this closure is disjoint from $X_{t_1}$. Hence
$(T_1 \to S', t'_1 \leadsto t_1, \xi_1)$
is an impurity of $\mathcal{F}_{S'}$ above $s'$.
In other words we have proved the contrapositive to part (2) of the
lemma. Finally, if $(T, t) \to (S, s)$ is an elementary
\'etale neighbourhood, then $(T_1, t_1) \to (S', s')$ is an
elementary \'etale neighbourhood too, and in this way we see that (1) holds.
\end{proof}

\end{document}
